\documentclass{melba}
% Submission: no option (line numbers on). arXiv during review: [arxiv]. Camera-ready: [accepted].

\usepackage{amsmath}

% The shared sections cite parenthetically as `Name~\cite{key}`; natbib's \cite is textual.
\let\cite\citep

\newcommand{\narrowfloatwidth}{\columnwidth}
\newcommand{\codeavailability}{Code, configurations, split manifests and the evaluation harness are
  available at \url{https://github.com/suxrobGM/morphoclip}.}

% Filled in by the publishing editor at camera-ready; not printed during review.
\melbaid{YYYY:NNN}
\doi{10.59275/j.melba.YYYY-AAAA}
\melbaauthors{Sukhrobbek Ilyosbekov, Shubham Gajjar and Rongfei Jin}
\email{ilyosbekov.s@northeastern.edu}
\volume{3}
\firstpageno{1}
\melbayear{2026}
\datesubmitted{yyyy-mm-dd}
\datepublished{yyyy-mm-dd}

\ShortHeadings{MorphoCLIP: Text-Supervised Perturbation Matching}{Ilyosbekov, Gajjar and Jin}

\title{MorphoCLIP: Text-Supervised Contrastive Learning for Perturbation Matching in Cell Painting
  Images}

\author{
  \name Sukhrobbek Ilyosbekov\aff{1},
  \name Shubham Gajjar\aff{1},
  \name Rongfei Jin\aff{1}
}
\affiliations{%
  \num 1 \addr Northeastern University, Portland, ME, USA
}

\abstract{%
  \input{sections/abstract}}

\keywords{Contrastive learning, Cell Painting, Morphological profiling, CLIP, Perturbation matching}

\begin{document}

\twocolumn[\maketitle]

\input{sections/introduction}
\input{sections/problem_statement}
\input{sections/related_work}
\input{sections/methods}
\input{sections/experiments}
\input{sections/results}
\input{sections/conclusion}

\acks{We thank the JUMP Cell Painting Consortium for making the CPJUMP1 data publicly available.
  AI writing assistants were used to edit the manuscript for language and clarity. All content,
  analyses and conclusions are the authors' own, and the authors take full responsibility for
  them.}

\ethics{The work follows appropriate ethical standards in conducting research and writing the
  manuscript. It uses only publicly available images of cultured human cell lines (U2OS and A549)
  from the CPJUMP1 dataset. No human subjects or animals were involved, so no institutional review
  board approval was required.}

\coi{We declare we don't have conflicts of interest.}

\data{CPJUMP1 is publicly available from the Cell Painting Gallery (\texttt{cpg0000-jump-pilot})
  on the Registry of Open Data on AWS. The reference metadata used by the benchmark, together with
  the code, configurations and split manifests, is available at
  \url{https://github.com/suxrobGM/morphoclip}.}

\bibliography{references}

\end{document}
